\documentclass[11pt]{article}
\usepackage[a4paper,margin=1in]{geometry}
\usepackage{booktabs}
\usepackage{longtable}
\usepackage{array}
\usepackage{hyperref}
\usepackage{xcolor}
\usepackage{fancyhdr}
\usepackage{enumitem}
\usepackage{caption}
\usepackage{amsmath}
\usepackage{graphicx}
\usepackage{titlesec}
\usepackage[english]{babel}

\hypersetup{colorlinks=true, linkcolor=blue!50!black, urlcolor=blue!50!black, citecolor=blue!50!black}
\titleformat{\section}{\normalfont\Large\bfseries}{\thesection}{1em}{}
\titleformat{\subsection}{\normalfont\large\bfseries}{\thesubsection}{1em}{}

\pagestyle{fancy}
\fancyhf{}
\fancyhead[L]{\small SIH26153 --- World Models for Predictive Cyber Defence}
\fancyhead[R]{\small\thepage}
\renewcommand{\headrulewidth}{0.3pt}

\newcommand{\good}{\textcolor{green!45!black}}
\newcommand{\bad}{\textcolor{red!55!black}}
\newcommand{\code}[1]{\texttt{#1}}

\title{\textbf{A World Model for Predictive Cyber Defence}\\
\large Learning Network State Dynamics to Forecast Attacker Progression}
\author{SIH26153 --- Technical Report}
\date{\today}

\begin{document}
\maketitle
\thispagestyle{empty}

\begin{abstract}
\noindent
This report documents the design, architecture search, and evaluation of a self-supervised
latent-transition model that learns from flow-level network telemetry and forecasts an
attacker's progression through a 7-stage MITRE ATT\&CK-derived taxonomy. The system is trained
on three public intrusion-detection datasets (CIC-IDS-2018, CIC-IDS-2017, CTU-13) pooled into
one 43-feature, 60-second-windowed representation, and evaluated with two supervised
downstream heads: (i) binary infiltration-probability forecasting 1--5 minutes ahead, and (ii)
7-way stage classification. \textbf{The result that should be read as primary is a cross-episode
holdout}: training on one of CIC-IDS-2018's two Infiltration episodes and testing purely on the
other (never mixed) gives AUC $\approx$ 0.68--0.82 --- a real, above-chance, but modest
generalization result, based on only the two such episodes the dataset contains. The commonly
quoted pooled/in-distribution numbers (AUC $\approx$ 0.97--0.98, obtained by training and
testing on different time-windows of the \emph{same} episodes) are reported for completeness
and cross-architecture comparison only and should \textbf{not} be read as a generalization
claim, since they can and likely do partly reflect episode-specific persistence rather than
transferable signal. Roughly a dozen architectural variants were compared, including a
Transformer-based model, a Recurrent State-Space Model (RSSM), domain-adversarial training, two
independent graph-structured encoders, and a KL-bottleneck ablation. The strongest,
seed-replicated configuration is an RSSM with an unfrozen encoder; unfreezing is shown here to
improve the two \emph{downstream classification/forecast heads} measured, which is a narrower
claim than an improvement in the underlying latent dynamics themselves --- this report's
evaluation is centered on those two task heads (Section~\ref{sec:evalscope} states this
boundary explicitly) and does not separately measure open-loop multi-step reconstruction error,
future communication-graph structure, or other direct measures of world-model quality. Several
stage classes have too few test occurrences to support a confident per-class claim (zero
Reconnaissance/Exfiltration occurrences, 5--17 Lateral-Movement/Initial-Access/Impact
occurrences at any forecast horizon); this is stated plainly rather than smoothed into the
headline numbers. The system includes a SHAP-based explanation of single-window stage
predictions (not of a full multi-step rollout --- a narrower claim, stated explicitly in
Section~\ref{sec:explain}) and a working live-inference prototype (Flask backend + browser
dashboard) accepting a CSV or raw PCAP file.
\end{abstract}

\tableofcontents
\newpage

%======================================================================
\section{Problem Framing}
%======================================================================

The problem statement (SIH26153) asks for a system that learns network state dynamics from
traffic telemetry, forecasts attacker progression, maps predictions onto MITRE ATT\&CK stages,
and exposes interpretable decision support --- explicitly favouring a \emph{world model}
(\(P(S_{t+1} \mid S_t)\), a model of how the network evolves) over a static classifier that only
labels the present. Two properties of the task shaped every subsequent design decision:

\begin{itemize}[leftmargin=1.5em]
  \item \textbf{Forecasting, not detection.} The deliverable is a probability of reaching a
  dangerous stage \(k\) steps into the future, not a same-timestep label. This rules out any
  architecture that only classifies the current window and pushes toward genuinely sequential
  models.
  \item \textbf{Interpretability is mandatory, not optional.} Both an attention/attribution
  channel (SHAP) and a closed-form, whiteboard-showable mechanism (a counted stage-transition
  matrix, \(\Pi\)) are required outputs, not an afterthought bolted on at the end.
\end{itemize}

%======================================================================
\section{Data and Feature Engineering}
%======================================================================

\subsection{Datasets}

\textbf{Provenance, stated precisely to avoid any ambiguity:} training used each dataset's
already-flow-level export --- CICFlowMeter's per-flow CSV output for CIC-IDS-2018/2017, and
Argus's per-flow \code{binetflow} export for CTU-13 --- re-aggregated by us from per-flow rows
into a common 60-second-window schema. \textbf{No raw packet capture (PCAP) was processed for
any of the three training datasets}; the public CIC-IDS-2018/2017 PCAP release is
$\sim$220\,GB and was not downloaded (Section~\ref{sec:limitations}). The only component of
this project that parses raw packets is the live-inference demo path
(Section~\ref{sec:demo}), which is separate from training and operates on user-supplied
PCAP/CSV input at inference time only.

\begin{itemize}[leftmargin=1.5em]
  \item \textbf{CIC-IDS-2018} (10 days, CICFlowMeter output) --- the primary source; carries
  Reconnaissance, Initial Access, Lateral Movement, Command \& Control, and Impact-stage
  attacks. No IP columns in the public release, which later blocked any host-communication
  graph representation for this dataset specifically.
  \item \textbf{CIC-IDS-2017} (8 days, \code{GeneratedLabelledFlows} release) --- folded into
  training on explicit instruction; retains real \code{Source IP}/\code{Destination IP} columns,
  unlike CIC-IDS-2018.
  \item \textbf{CTU-13} (13 scenarios, Argus binetflow) --- 13 independent real botnet captures
  (Neris, Rbot, Virut, Murlo, Sogou, and others), contributing mainly Benign/Command-and-Control
  labels; also carries real \code{SrcAddr}/\code{DstAddr} and real TTL.
\end{itemize}

\noindent Two datasets suggested by the problem statement were evaluated and \textbf{not}
integrated, for stated, verified reasons (Section~\ref{sec:limitations}): CICIoT2023 (no
timestamps, packets pre-shuffled into fixed-size groups --- structurally incompatible with
sequence forecasting) and raw PCAP-level capture for CIC-IDS-2018/2017 (the public PCAP release
is $\sim$220\,GB, judged not worth the bandwidth/storage cost for the marginal features gained
over the flow-level CSVs).

\subsection{Stage-label mapping: provenance and scope}
\label{sec:attackmapping}

None of the three datasets ships native MITRE ATT\&CK labels. The 7-stage taxonomy used
throughout this report (\code{Benign}, \code{Reconnaissance}, \code{Initial Access},
\code{Lateral Movement}, \code{Command and Control}, \code{Exfiltration}, \code{Impact}) is
obtained by keyword-matching each dataset's own native attack-type label string (e.g.\
\code{"DDoS attacks-LOIC-HTTP"}, \code{"Infilteration"}, \code{"PortScan"}) against a fixed,
first-match-wins rule table (reproduced in full below), then assigning that single stage label
to every flow in the window and taking the majority label as the window's stage. This mapping
is a \textbf{first-class, reviewable modelling decision, not native ground truth}, and two
consequences follow directly from that:

\begin{itemize}[leftmargin=1.5em]
  \item \textbf{The 7 stages are treated as one label per window, in a fixed evaluation order}
  --- real MITRE ATT\&CK techniques and tactics can co-occur, repeat, and branch, and this
  project does not claim otherwise. The single-label-per-window taxonomy is an evaluation
  simplification chosen because the source datasets carry a single attack-type string per flow
  in the first place; it is not a modelling claim that real attacker behaviour is linear.
  \item \textbf{Two mapping decisions are judgment calls, stated here rather than left implicit
  in code}: CIC-IDS-2018's \code{"Infilteration"} label (its multi-phase exploit-delivery-plus
  internal-portscan scenario) is mapped to \emph{Lateral Movement} rather than \emph{Initial
  Access}, on the basis that the labeled traffic is the attacker already inside and moving
  through the network; and CTU-13's \code{"Background"} traffic (unconfirmed ambient traffic,
  not attested as either attack or known-good) is mapped to \emph{Benign} because it is not
  confirmed malicious, not because it is confirmed clean.
\end{itemize}

\noindent The full first-match-wins keyword table:

\begin{center}
\small
\begin{tabular}{ll}
\toprule
\textbf{Keyword (case-insensitive substring)} & \textbf{Stage} \\
\midrule
benign, normal, background & Benign \\
portscan, port scan, recon & Reconnaissance \\
bruteforce, brute force, patator, sql injection, xss, web attack & Initial Access \\
infilt, heartbleed & Lateral Movement \\
bot, botnet & Command and Control \\
exfil & Exfiltration \\
ddos, dos(-/\_/space), goldeneye, slowloris, slowhttptest, hulk & Impact \\
\bottomrule
\end{tabular}
\end{center}

\noindent An unrecognized non-benign label raises an error at preprocessing time rather than
silently defaulting to Benign (\code{attack\_map.py}'s \code{stage\_idx\_for\_label}) --- every
label that entered training is accounted for in this table, none were dropped silently. The
$\Pi$ transition matrix (Section~\ref{sec:pi_context}) is counted directly over this same
per-window single-stage label sequence; it should be read as a baseline/explanation aid over
that simplified label stream, \textbf{not as a claim about ATT\&CK's real non-linear
structure}.

\subsection{Feature schema}

Each 60-second window is reduced to a flat 43-dimensional vector: flow volume/rate statistics,
TCP flag rates, protocol shares, destination-port structure (entropy, well-known-service shares,
distinct-port fraction --- the scan signature), connection-shape signals (one-way/zero-byte
fractions), and packet-level slots (TTL, fragmentation, retransmission, payload-size variance)
--- populated with real values where the source dataset supports it (CTU-13's real TTL) and
zero-filled with an explicit flag otherwise (CIC-IDS-2018/2017's flow-CSV-only public release).
No IP or port identity is ever fed to the model as a feature, only used as an aggregation key ---
a deliberate guard against the model memorizing ``this IP is the attacker'' instead of learning a
transferable traffic signature.

\subsubsection*{The log1p fix}
Heavy-tailed count/volume columns (flow count, byte/packet means) span orders of magnitude ---
a benign window might see $n\_flows \approx 1{,}300$ while a DDoS window reaches $n\_flows
\approx 275{,}000$. Naive robust (median/IQR) scaling alone clips anything past $\pm 8$
IQR-units to one identical saturated value, which erases exactly the signal that should most
clearly mark an Impact-stage window. Applying $\log(1+x)$ before scaling compresses this range
without destroying relative ordering. This was a concrete, verified bug fix, not a hypothetical
concern: Impact-stage detection was materially better after the fix.

%======================================================================
\section{Core Architecture}
%======================================================================

\subsection{Two-phase design}

\begin{enumerate}[leftmargin=1.5em]
  \item \textbf{Phase 1 --- self-supervised dynamics.} An encoder maps the 43-dim window to a
  latent $z_t$; a transition model produces a distribution $\mathcal{N}(\mu, \sigma^2)$ over
  $z_{t+1}$ (never a point estimate); a decoder reconstructs $\hat S_{t+1}$. Trained purely on
  Gaussian negative log-likelihood of the real next window, with \textbf{no label ever touched}.
  \item \textbf{Phase 2 --- supervised heads.} With the Phase-1 backbone optionally frozen
  (Section~\ref{sec:unfreezing} covers when it should not be), a linear \code{StageHead}
  (7-way MITRE classification) and \code{InfilHead} (binary infiltration probability) are
  trained on top, plus a Laplace-smoothed, counted stage-transition matrix $\Pi$\label{sec:pi_context}
  (never gradient-trained, and counted over the simplified single-stage-per-window label stream
  described in Section~\ref{sec:attackmapping} --- see that section for why this is a
  baseline/explanation aid, not a claim about ATT\&CK's real structure) for a cheap,
  interpretable closed-form forecast check: $\hat\gamma_{t+K} = (\Pi^\top)^K \gamma_t$.
\end{enumerate}

\subsection{The multi-step training fix}\label{sec:multistep}

The dynamics loss is scored two ways, \textbf{concatenated} into one loss rather than weighted
against each other:
\begin{itemize}[leftmargin=1.5em]
  \item one-step, teacher-forced: real $z_{\le t}$ in, score $z_{t+1}$ against reality, at every
  position;
  \item multi-step, open-loop: encode only the context, roll the transition forward
  \emph{consuming nothing but its own sampled output}, decode each imagined step, score against
  the real future observations.
\end{itemize}
Training on the one-step term alone never teaches the model to stay accurate over an open-loop
rollout --- it only ever sees real history during training, then is asked to imagine chains of
its own predictions at inference. Adding the multi-step term (this is what ``latent overshooting''
means in the RSSM literature) closed nearly the entire real-vs-imagined accuracy gap for the two
best-supported classes (Benign, Command \& Control), while correctly leaving the three
data-scarce classes untouched --- confirming this was a training-regime fix, not a data-volume
fix.

\subsection{Infiltration probability is derived, not a dedicated head}

Infiltration probability is the stage-classifier's own softmax mass on
$\{\text{Lateral Movement}, \text{Command \& Control}, \text{Exfiltration}\}$, not a
separately-trained head. A dedicated \code{InfilHead} was built and tested twice, under two
different class-weighting schemes, and was \textbf{consistently worse} than the derived
approach both times --- a clean, confirmed negative result, reported honestly rather than
buried.

\subsection{Evaluation scope: what is, and is not, measured}
\label{sec:evalscope}

Every result in Sections~\ref{sec:archsearch}--\ref{sec:pstage} is a measurement of the
\textbf{two downstream heads} (\code{StageHead}, \code{InfilHead}) built on top of the Phase-1
latent, not of the Phase-1 dynamics model in isolation. This distinction matters because a
change that improves those heads --- most notably unfreezing the encoder
(Section~\ref{sec:unfreezing}) --- is not automatically evidence that the underlying
\emph{learned network dynamics} improved; it may equally mean the representation was pulled
toward whatever the classification/forecasting objective rewards, at some unmeasured cost to
the dynamics model's own quality. Three concrete gaps follow from this, stated explicitly
rather than left implicit:

\begin{itemize}[leftmargin=1.5em]
  \item \textbf{No separate open-loop multi-step reconstruction metric is reported.} The
  multi-step training loss (Section~\ref{sec:multistep}) trains against exactly this
  quantity --- Gaussian NLL of the decoder's imagined rollout against the real future window
  --- but this report surfaces only the downstream Brier/AUC/F1 built on top of that latent,
  never the raw per-feature reconstruction error itself, frozen vs.\ unfrozen, side by side.
  \item \textbf{No future communication-structure/graph-level prediction is evaluated.} The
  two graph-structured encoder variants (Section~\ref{sec:graphs}) are evaluated only through
  the same two downstream heads as every other variant; neither this report nor the underlying
  project measures whether either graph encoder's own predicted future edge structure matches
  the real future communication pattern.
  \item \textbf{SHAP explains one window's classification, not a multi-step forecast rollout}
  (Section~\ref{sec:explain} states this boundary explicitly).
\end{itemize}

\noindent The system therefore supports a \textbf{narrower, more defensible claim} than
``the world model's dynamics improved'': it supports ``the latent representation, combined
with two supervised heads, produces a useful forecast and classification signal, whose
quality is measured directly; whether the \emph{unsupervised dynamics model in isolation}
improved is not separately established here.'' Readers should treat every comparison in
Section~\ref{sec:archsearch} under that scope.

%======================================================================
\section{Architecture Search}
\label{sec:archsearch}
%======================================================================

Roughly a dozen variants were built and evaluated on identical splits, using two metrics
throughout: \textbf{Brier score} (mean squared error of the predicted probability against the
binary outcome --- calibration) and \textbf{AUC} (probability a random true positive outranks a
random true negative --- ranking quality, ignoring calibration; 0.5 is random, below 0.5 is
\emph{worse} than random). Table~\ref{tab:master} summarizes every architecture on pooled
metrics; per-domain and per-class detail follows.

\begin{table}[htbp]
\centering
\caption{Every architecture tried, pooled (``all domains'') metrics, neural method.}
\label{tab:master}
\small
\begin{tabular}{lccc}
\toprule
\textbf{Architecture} & \textbf{Brier ($k$1$\to$5)} & \textbf{AUC ($k$1$\to$5)} & \textbf{macroF1 (real/imag.)} \\
\midrule
Baseline Transformer (frozen)        & 0.140 $\to$ 0.128 & 0.913 $\to$ 0.911 & 0.376 / 0.302 \\
RSSM, plain (frozen)                 & 0.169 $\to$ 0.174 & 0.931 $\to$ 0.909 & 0.329 / 0.273 \\
RSSM + domain-adversarial (frozen)   & 0.188 $\to$ 0.188 & 0.927 $\to$ 0.912 & 0.295 / 0.327 \\
RSSM + feature-graph encoder (frozen)& 0.124 $\to$ 0.141 & 0.955 $\to$ 0.931 & 0.386 / 0.306 \\
KL-bottleneck-on-Transformer ablation& 0.126 $\to$ 0.098 & 0.935 $\to$ 0.939 & 0.205 / 0.337 \\
\textbf{RSSM, encoder unfrozen}       & \textbf{0.054 $\to$ 0.061} & \textbf{0.982 $\to$ 0.968} & \textbf{0.521 / 0.467} \\
RSSM unfrozen + DANN (seed 1)         & 0.048 $\to$ 0.058 & 0.983 $\to$ 0.977 & 0.519 / 0.556 \\
RSSM unfrozen + DANN (seed 2)         & 0.050 $\to$ 0.052 & 0.976 $\to$ 0.967 & 0.671 / 0.437 \\
G-RSSM (host-graph, 2 of 3 domains)   & 0.174 $\to$ 0.181 & 0.775 $\to$ 0.765 & 0.418 / 0.469$^\dagger$ \\
\bottomrule
\end{tabular}
\vspace{2pt}
\\ \raggedright\footnotesize $^\dagger$ Evaluated on a 3-of-7-class test set (only CTU-13 and
CIC-IDS-2017, which lack several stages in their test slices) --- not directly comparable to the
6-class pooled numbers above.
\end{table}

\subsection{Domain-adversarial training (DANN)}

Pooling three genuinely different network environments causes a real, measured failure: without
correction, the stage classifier and $\Pi$ work well on whichever domain dominates training and
are actively worse than random on the others. DANN adds a small domain classifier trained
\emph{adversarially} against the shared encoder via a Gradient Reversal Layer (GRL): the domain
head itself gets better at guessing which dataset a window came from, while the encoder
underneath receives the \emph{negated} gradient, pushed to make that guess harder. On the plain
Transformer this helped the worst domain (CIC-IDS-2017) but hurt the previously-best one
(CTU-13) --- a genuine trade-off, not a clean win.

\subsection{Two independent graph-structured encoders}\label{sec:graphs}

Two structurally different graph ideas were tried, with markedly different outcomes:

\begin{description}[leftmargin=1.5em,style=nextline]
  \item[Feature-level graph (GraFT-inspired).] Nodes are the 43 \emph{features themselves}
  (not network hosts); edges are (a) fixed domain-knowledge groupings (e.g.\ all TCP-flag
  features connected to each other) and (b) per-window dynamic top-$k$ cosine-similarity edges.
  A relational GCN with per-relation-type weights replaces the flat MLP encoder. Combined with
  the Transformer, this was \textbf{unstable} across epoch counts (AUC swinging from
  $\approx$0.5 to $\approx$0.07 between 30- and 80-epoch runs on CIC-IDS-2017) --- consistent
  with the literature on GNN oversmoothing/initialization sensitivity, and reverted. Combined
  with RSSM instead, it was one of the \textbf{best} results of the whole project (best pooled
  AUC of any frozen variant, best frozen-model classification F1) --- the instability was tied
  to the Transformer's own attention-based transition competing with a second attention
  mechanism inside the encoder, not an inherent flaw of the graph idea.
  \item[Host-communication graph (G-RSSM, arXiv:2604.14811).] Nodes are actual network hosts
  (IP addresses); edges are observed communication within a 60-second window; a per-host GRU
  plus masked cross-host attention replaces the flat encoder entirely. Only CTU-13 and
  CIC-IDS-2017 carry host identity (CIC-IDS-2018's public CSVs do not), so this variant covers
  two of three domains. The first implementation produced forecasts \emph{systematically
  anti-correlated} with ground truth (AUC 0.13--0.30, below random) --- traced to a real,
  identifiable bug: the top-16 hosts selected for each window were re-chosen \emph{independently
  per window}, so a given ``node slot'' silently referred to a different physical host from one
  minute to the next, making the per-host recurrence meaningless. Fixing this (a day-level
  fixed host roster, one consistent slot per host for the whole day) genuinely repaired
  CIC-IDS-2017 forecasting (AUC 0.14 $\to$ 0.58, $\Pi$-check reaching 0.85) but left CTU-13 still
  broken (AUC 0.20--0.31) --- a second, distinct, unresolved issue (likely that CTU-13's
  low-volume botnet hosts never make a naive top-16-by-volume cut).
\end{description}

\subsection{The KL-bottleneck ablation: why does RSSM beat the Transformer?}

RSSM's advantage over the Transformer could come from at least two different mechanisms: its
GRU recurrence/smaller parameter count, or its stochastic posterior/prior plus free-bits KL
regularization (the ``bottleneck''). To isolate this, a small posterior head was added
\emph{on top of the otherwise-unchanged Transformer} --- same backbone, same parameter count, no
GRU --- trained with DreamerV3's exact balanced free-bits KL against the Transformer's own
prior. Result: the KL-bottleneck alone improved pooled Brier/AUC over the plain Transformer and
won outright on CTU-13, but did \textbf{not} reproduce RSSM's CIC-IDS-2017 strength --- so the
GRU/recurrence itself (or seed variance) explains that specific domain's result, not the KL
mechanism alone. The hypothesis splits cleanly rather than resolving to one cause.

\subsection{Encoder freezing: the real bottleneck for classification}
\label{sec:unfreezing}

Every architecture above trains its \code{StageHead} on a \textbf{frozen} Phase-1 encoder ---
shaped only by the label-blind dynamics loss. A raw-feature logistic regression proved the seven
stages are linearly separable in the 43-dim input space (one-vs-rest AUC 0.90--0.999, including
Reconnaissance with only 12 training examples, and Initial Access at 0.999), yet every frozen
architecture scored $\le$0.402 F1 on Initial Access and $0.000$ on it for most variants. The
mechanism: a frozen, label-blind encoder has no incentive to preserve a rare event's
distinguishing detail if it barely moves the averaged dynamics loss --- exactly the
rate-distortion trade-off that makes rare, ``surprising'' states the most expensive thing a
KL-regularized latent can encode faithfully.

\textbf{Unfreezing the encoder} for Phase 2 (letting the classification/infiltration gradient
flow back into it) confirmed this decisively: macroF1 jumped $0.329 \to 0.521$, infiltration
Brier improved roughly 3$\times$, Initial Access F1 went from \bad{0.000} to \good{0.402}, Impact
from \bad{0.208} to \good{0.628}. The cost: RSSM's own CIC-IDS-2017 domain-generalization
strength partially collapsed (AUC $0.87 \to 0.65$), because fine-tuning on labels pulls the
representation toward whatever the label-heavy domains need. Combining unfreezing with DANN
recovered most of that loss on one seed (AUC up to 0.83) but the recovery \textbf{did not
replicate} on a second seed (AUC 0.19--0.30) --- see Section~\ref{sec:seedvariance}.

\subsection{Seed variance: the discipline that mattered most}
\label{sec:seedvariance}

RSSM showed large run-to-run variance throughout this project. Two seeds of the same
``unfrozen + DANN'' configuration gave opposite answers to the specific question ``does DANN
recover RSSM's lost CIC-IDS-2017 strength'': seed 1 said yes (AUC $0.73\to0.83$), seed 2 said no
(AUC $0.19\to0.30$, worse than a coin flip at $k{=}1$). Everything else about that configuration
--- pooled metrics, CIC-IDS-2018, CTU-13, classification F1 --- replicated closely across both
seeds. The honest conclusion: unfreezing and the general RSSM architecture are solid, replicated
findings; the specific claim that DANN reliably rescues CIC-IDS-2017 domain generalization is
\textbf{not established} and should not be reported as a settled result.

%======================================================================
\section{Per-Stage Forecasting Results}\label{sec:pstage}
%======================================================================

Binary infiltration forecasting collapses the full 7-way stage distribution into one number
(probability mass on the three post-compromise stages) because it is the decision-relevant
question for an operator (``will this become a serious breach'') and because several stages have
too few test occurrences for a trustworthy per-horizon AUC. Table~\ref{tab:perstage} makes this
concrete, using \code{RSSM, encoder unfrozen} pooled across all three domains.

\begin{table}[htbp]
\centering
\caption{Per-stage forecasting, pooled, \code{derived}/\code{pi} AUC at $k{=}1$ and $k{=}5$.}
\label{tab:perstage}
\small
\begin{tabular}{lccc}
\toprule
\textbf{Stage} & \textbf{$n$ positive ($k$1$\to$5)} & \textbf{AUC $k{=}1$} & \textbf{AUC $k{=}5$} \\
\midrule
Benign                & 711 $\to$ 731 & 0.960 & 0.951 \\
Command and Control   & 989 $\to$ 977 & \textbf{0.987} & \textbf{0.974} \\
Initial Access        & 17 $\to$ 13 (small) & 0.979 & 0.999 \\
Lateral Movement      & 14 $\to$ 10 (small) & 0.947 & 0.903 \\
Impact                & 5 $\to$ 5 (not trustworthy) & 1.000 & 1.000 \\
Reconnaissance        & 0 everywhere & \multicolumn{2}{c}{insufficient test occurrences} \\
Exfiltration          & 0 everywhere & \multicolumn{2}{c}{insufficient test occurrences} \\
\bottomrule
\end{tabular}
\end{table}

Command-and-Control forecasting is the one genuinely large-sample, statistically solid result
across every domain. Reconnaissance and Exfiltration cannot be forecast at all with the current
data: there are zero test windows anywhere that transition into either stage within any 5-minute
horizon --- a data-scarcity floor, not an architecture limitation.

%======================================================================
\section{Baselines and the Cross-Episode Generalization Test}\label{sec:crossepisode}
%======================================================================

\subsection{The logistic-regression baseline}

A per-horizon logistic regression (five separate binary models, one per forecast horizon,
trained on the flattened context window, class-weight balanced) was built as the mandated
baseline and given every fair chance. For most of the architecture search, it was competitive
with or dominant over the world model --- a result reported honestly throughout rather than
minimized, and the single strongest signal that architectural complexity was not obviously
buying its keep until the final \code{RSSM unfrozen} configuration.

\subsection{The memorization check}

CIC-IDS-2018 contains exactly two Infiltration episodes (Wednesday-28-02-2018 and
Thursday-01-03-2018, both mapped to Lateral Movement). Every LR number reported elsewhere in
this project was obtained by pooling windows from \emph{all} days (including both infiltration
days) and testing on a later time slice of those same days --- meaning the model had already
partially seen the episode it was evaluated on. A true \textbf{episode-level} holdout (train the
LR purely on one day, test purely on the other, never mixed) tells a very different story:

\begin{table}[htbp]
\centering
\caption{True cross-episode holdout, logistic regression.}
\small
\begin{tabular}{lcc}
\toprule
\textbf{Direction} & \textbf{Brier ($k$1$\to$5)} & \textbf{AUC ($k$1$\to$5)} \\
\midrule
Train Wed-28 \(\to\) Test Thu-01 & 0.182 $\to$ 0.220 & 0.792 $\to$ 0.824 \\
Train Thu-01 \(\to\) Test Wed-28 & \bad{0.584 $\to$ 0.594} & 0.676 $\to$ 0.742 \\
\bottomrule
\end{tabular}
\end{table}

\noindent AUC collapses from the reported $\approx$0.995 (within-episode) to 0.68--0.82 (true
holdout) --- a substantial drop, though not to chance, meaning LR does learn some genuinely
transferable signal, just far less than the in-distribution number suggested. The
Thursday$\to$Wednesday direction is worse than \emph{useless} on Brier ($0.58$, worse than
predicting the constant base rate), even though its AUC ranking is still reasonably above random
--- a real dissociation between ranking quality and calibration trustworthiness.

The same test on the world model (\code{RSSM unfrozen}, trained on everything \emph{except} the
held-out day) shows AUC 0.72--0.82 across both directions and, critically, \textbf{never} the
catastrophic Brier failure LR showed in one direction (Brier stays 0.21--0.31 both ways). Neither
method has a decisive AUC edge on a genuinely unseen episode; the real differentiator is
calibration stability, not ranking accuracy --- a materially more honest and more useful claim
for a predictive-defence system than ``the model beats the baseline.''

%======================================================================
\section{Explainability}\label{sec:explain}
%======================================================================

\subsection{Attention and $\Pi$}

Attention weights are read directly off the Transformer's own computation, not a post-hoc
approximation, when that backbone is used --- but reading an attention weight off a model is
not, on its own, a guarantee that the weight is a \emph{faithful} explanation of the model's
actual computation (a well-documented general limitation of attention-as-explanation in the
interpretability literature); it is offered here as a diagnostic signal, not a certified
causal account. The $\Pi$ transition matrix is directly interpretable by construction ---
every entry is a literal, counted transition probability, and the $K$-step forecast is a
matrix power, not a learned function --- subject to the scope stated in
Section~\ref{sec:attackmapping} (it is counted over the simplified single-stage-per-window
label, not native ATT\&CK ground truth).

\subsection{SHAP}

A \code{KernelExplainer} wraps the model's stage-classification output (holding the fixed prior
context and perturbing only the most recent window's 43 raw features) and explains, per window,
which features drove the \textbf{current, single-window} stage call. \textbf{This is explicitly
not the same claim as explaining a multi-step ($k{=}1..5$) forecast rollout} --- no SHAP
attribution is computed against the rollout's own output, and none is claimed to exist here;
extending this to a genuine multi-step-rollout explanation is a stated open item, not an
implicit claim of this report. Two consistent, load-bearing findings about the single-window
explanation that \emph{is} computed:

\begin{itemize}[leftmargin=1.5em]
  \item \textbf{SHAP magnitude tracks decision-boundary proximity, not a fixed scale.} On
  confidently-Benign windows, attributions are tiny ($\sim$0.0001--0.0005) because the softmax
  has little room to move; during a real, live-tested Lateral Movement attack, attributions grow
  roughly 30--100$\times$ larger ($\sim$0.004--0.016), correctly reflecting that the model is
  actually near a decision boundary.
  \item \textbf{Validated against a genuine attack in real time.} Replaying the actual
  Wednesday-28-02-2018 attack window-by-window (Section~\ref{sec:demo}) showed infiltration
  probability climbing \emph{one window before} the ground-truth label flipped to Lateral
  Movement (0.61--0.67 at the window immediately prior), then the stage classification itself
  correctly flipping to Lateral Movement at the true transition window and staying correct for
  the sustained 18-window attack duration, with SHAP consistently attributing the call to
  \code{svc\_smtp\_frac}, \code{svc\_ssh\_frac}, \code{half\_open\_frac}, and
  \code{bwd\_pkts\_mean}.
\end{itemize}

%======================================================================
\section{Working Prototype}
\label{sec:demo}
%======================================================================

Two deliverables were built beyond the research pipeline:

\begin{description}[leftmargin=1.5em,style=nextline]
  \item[A replay dashboard.] Visualizes the model's \emph{saved} predictions over a real attack
  day as an animated MITRE-stage strip, an infiltration-probability forecast chart, a scrubbable
  timeline, and per-window SHAP bars. Useful for demonstration and inspection, but --- stated
  plainly in its own footer --- it is a report of precomputed output, not a live tool.
  \item[A live-inference server (the actual PS deliverable).] A Flask backend (\code{server.py})
  loads the trained model once at startup and exposes two endpoints: \code{POST /api/analyze}
  accepts an uploaded CSV (43-feature schema) or a raw PCAP file, extracts windows (with genuine
  packet-level TTL/fragmentation/payload-size features when given a PCAP, computed from real
  packets rather than zero-filled), and runs the world model \emph{live} to return a forecast and
  stage classification for every window; \code{POST /api/explain} computes a SHAP explanation for
  one specific window on demand, so a full KernelExplainer pass is never paid for windows the
  user never looks at. A browser dashboard (drag-and-drop upload, or one-click real sample
  buttons for the two genuine CIC-IDS-2018 attack days) drives this API and renders the same
  visual language as the replay dashboard, but every number on screen was computed by the model
  for that specific input at request time.
\end{description}

%======================================================================
\section{Known Limitations}
\label{sec:limitations}
%======================================================================

Reported plainly, as an inventory of open problems rather than a list of excuses:

\begin{itemize}[leftmargin=1.5em]
  \item \textbf{Reconnaissance and Exfiltration are unforecastable and (Reconnaissance) barely
  classifiable with current data} --- 12 total training examples for Reconnaissance, zero
  Exfiltration examples anywhere in the pooled corpus. No architecture change fixes a class the
  data does not contain.
  \item \textbf{CTU-13's host-graph forecasting (G-RSSM) remains broken} after fixing the
  node-identity bug; a second, distinct issue (likely host-selection criterion) is unresolved.
  \item \textbf{CICIoT2023 could not be integrated}: its public release has no timestamps and
  pre-shuffles packets into fixed-size groups, incompatible with time-windowed sequence
  forecasting at the source-format level, not a preprocessing choice.
  \item \textbf{The strongest single result (RSSM unfrozen + DANN) has an unreplicated headline
  claim.} Its domain-generalization recovery on CIC-IDS-2017 held on one seed and failed on a
  second; only the plain \code{RSSM, encoder unfrozen} configuration (without DANN) has fully
  replicated numbers across the seeds actually run.
  \item \textbf{Packet-level PCAPs for CIC-IDS-2018/2017 were not downloaded} ($\sim$220\,GB);
  packet-level features (TTL, fragmentation, retransmission, payload-size variance) are
  zero-filled for those two datasets in the training pipeline.
  \item \textbf{The live-PCAP inference path has an unmeasured train-serving distribution
  mismatch.} When the live demo is given a real PCAP, those same packet-level fields are
  populated with genuine, non-zero values computed from real packets --- values the model
  never observed during training on the two zero-filled datasets (roughly two-thirds of the
  pooled training corpus by day count). Whether this shift degrades reliability has not been
  measured (e.g.\ via a held-out comparison of the model's calibration on zero-filled vs.\
  genuinely-populated packet-level inputs); the live demo's few qualitative test runs are not
  a substitute for that measurement, and it should be treated as open rather than assumed
  benign.
\end{itemize}

%======================================================================
\section{Conclusion}
%======================================================================

The system addresses the problem statement's core requirements, within the scope this report
states explicitly rather than implies: a latent-transition model (not a static classifier)
built from a self-supervised dynamics objective supports downstream forecasting of attacker
progression several minutes ahead and 7-way stage classification, maps predictions onto a
MITRE-ATT\&CK-derived taxonomy whose construction is documented in
Section~\ref{sec:attackmapping} rather than left implicit, and exposes two explanation
channels (attention/$\Pi$ and single-window SHAP) whose claim boundaries are stated in
Section~\ref{sec:explain}. The one result that should be treated as the project's actual
generalization claim is the cross-episode holdout (Section~\ref{sec:crossepisode}, AUC
$\approx$0.68--0.82 on two held-out episodes) --- modest, above chance, and honestly small-sample.
The strongest configuration on the (separately reported, in-distribution) pooled metrics,
\code{RSSM with an unfrozen encoder}, replicates across five seeds on those metrics, but that
result describes the two downstream heads measured, not an independently verified improvement
in the dynamics model itself (Section~\ref{sec:evalscope}). Equally important to the project's
credibility is what is reported \emph{without} spin: a plain logistic regression was
competitive for most of the search, a dedicated infiltration head made things worse, two
graph-encoder attempts each failed in different, diagnosable ways, one headline result
(DANN's domain-generalization recovery) did not survive a second seed, several stage classes
have too little test support for a confident claim, and the true cross-episode generalization
ceiling is substantially lower than the in-distribution numbers on their own would imply.
Every one of these negative or
qualified findings was investigated to a specific, verifiable cause rather than smoothed over
--- which is, together with the working live-inference prototype, the actual contribution of
this project.

%======================================================================
\section*{Reproducibility}
%======================================================================

Provided here, stated honestly against what is \textbf{not} provided (this report is not a
sealed reproducibility package):

\begin{itemize}[leftmargin=1.5em]
  \item \textbf{Provided:} context length 16 windows, forecast horizon 5 windows; Phase~1
  (dynamics) trained 30 epochs, Phase~2 (heads) 20 epochs, AdamW, weight decay $10^{-4}$,
  learning rate $3\times10^{-4}$ (Phase~1) / $10^{-3}$ (Phase~2), gradient clipped to norm 1.0;
  RSSM dimensions $d_h{=}32$, $d_z{=}16$, $d_{\mathrm{embed}}{=}64$; free-bits KL floor of 1
  nat per direction (DreamerV3's exact balanced free-bits formulation); temporal
  train/val/test split 70/15/15 per source file, never crossing a day/scenario boundary; the
  five-seed replication of \code{RSSM, encoder unfrozen} reported in
  Section~\ref{sec:archsearch} used five independent Phase-1+Phase-2 training runs with no
  fixed random seed (relying on run-to-run variance itself, not a chosen seed list) and its
  full per-seed metrics.
  \item \textbf{Not provided:} exact train/val/test split manifests (row-level file lists),
  checkpoint file hashes, a frozen/sealed test protocol, or dataset-preprocessing content
  hashes. Reproducing the headline table from this report alone would require re-running the
  described pipeline against the same public dataset releases and accepting the same run-to-run
  variance documented in Section~\ref{sec:archsearch}, not a byte-identical replay.
\end{itemize}

%======================================================================
\section*{References}
%======================================================================

\begin{enumerate}[leftmargin=1.5em]
  \item Hafner, D., Pasukonis, J., Ba, J., \& Lillicrap, T. (2023). \emph{Mastering Diverse
  Domains through World Models.} arXiv:2301.04104. (DreamerV3; source of the exact balanced
  free-bits KL formulation used throughout this project's RSSM variants.)
  \item Ganin, Y., \& Lempitsky, V. (2015). \emph{Unsupervised Domain Adaptation by
  Backpropagation.} Proceedings of ICML 2015. (Gradient Reversal Layer / domain-adversarial
  training, Section~\ref{sec:archsearch}.)
  \item Schlichtkrull, M., Kipf, T. N., Bloem, P., van den Berg, R., Titov, I., \& Welling, M.
  (2018). \emph{Modeling Relational Data with Graph Convolutional Networks.} ESWC 2018.
  (Relational GCN formulation underlying the feature-level graph encoder,
  Section~\ref{sec:graphs}.)
  \item AAAI'26 GraFT paper (\url{https://ojs.aaai.org/index.php/AAAI/article/view/40029}) ---
  motivating reference for the feature-level relational graph encoder design.
  \item Karacelebi, C., Sahin, Y. T., Surer, E., \& Onur, E. \emph{Learning Ad Hoc Network
  Dynamics via Graph-Structured World Models.} arXiv:2604.14811. (G-RSSM; per-host
  latent-state, cross-host-attention design underlying the host-communication graph encoder,
  Section~\ref{sec:graphs}.)
  \item Sharafaldin, I., Lashkari, A. H., \& Ghorbani, A. A. (2018). \emph{Toward Generating a
  New Intrusion Detection Dataset and Intrusion Traffic Characterization.} ICISSP 2018.
  (CIC-IDS-2018/2017, Canadian Institute for Cybersecurity.)
  \item Garcia, S., Grill, M., Stiborek, J., \& Zunino, A. (2014). \emph{An Empirical
  Comparison of Botnet Detection Methods.} Computers \& Security, 45, 100--123. (CTU-13,
  Stratosphere IPS Laboratory.)
\end{enumerate}

\end{document}
